% === Appendix ===
\newpage
\appendix
\setcounter{table}{0}
\setcounter{figure}{0}
\setcounter{section}{0}
\renewcommand{\thetable}{\Alph{section}\arabic{table}}
\renewcommand{\thefigure}{\Alph{section}\arabic{figure}}
\renewcommand{\thesection}{\Alph{section}}

% --- Appendix Table of Contents (toggle via \appendixtoctrue / \appendixtocfalse in main.tex) ---
\ifappendixtoc
  \startcontents[appendix]
  \noindent{\LARGE\bfseries Appendix}
  \vspace{1.5em}
  \titlecontents{lsection}[0em]{\large\bfseries\color{msftblue}\vspace{0.8em}}{\thecontentslabel\hspace{0.5em}}{}{\titlerule*[0.5pc]{.}\contentspage}
  \titlecontents{lsubsection}[2.5em]{\normalsize\color{msftblue}\vspace{0.3em}}{\thecontentslabel\hspace{0.5em}}{}{\titlerule*[0.5pc]{.}\contentspage}
  \printcontents[appendix]{l}{1}{\setcounter{tocdepth}{2}}
  \newpage
\fi

% ============================================================
%  Appendix A: Additional Results
% ============================================================
\section{Additional Results}
\label{sec:additional_results}

This section provides supplementary experimental results that could not be included in the main paper due to space constraints.

\subsection{Extended Comparison}

% --- Example: additional results table in appendix ---
\begin{table}[h!]
  \centering
  \caption{Extended comparison with additional baselines on [your benchmark].}
  \label{tab:extended_results}
  \begin{tabular}{lccc}
    \toprule
    Method & Metric A ($\uparrow$) & Metric B ($\downarrow$) & Metric C ($\uparrow$) \\
    \midrule
    Baseline D & 72.1 & 0.48 & 65.3 \\
    Baseline E & 74.8 & 0.44 & 68.7 \\
    Baseline F & 77.3 & 0.39 & 71.2 \\
    \midrule
    \method (Ours) & \textbf{82.1} & \textbf{0.31} & \textbf{76.9} \\
    \bottomrule
  \end{tabular}
\end{table}

\subsection{Additional Visualizations}

% --- Example: appendix figure ---
\begin{figure}[h!]
  \centering
  \fbox{\rule{0pt}{1.5in}\rule{0.9\linewidth}{0pt}}
  \caption{Additional qualitative results. Replace with your supplementary visualizations.}
  \label{fig:additional_vis}
\end{figure}

\subsection{Implementation Details}

Provide additional implementation details here, such as hyperparameter settings, data preprocessing pipelines, or hardware specifications that are too detailed for the main text.

% ============================================================
%  Appendix B: Prompt Templates
% ============================================================
\section{Prompt Templates}
\label{sec:prompt_templates}

This section documents the prompts used in our pipeline.

\subsection{Evaluation Prompt}

% --- Example: system prompt with natural language ---
\begin{promptbox}[System Prompt for Evaluation]
\noindent\textbf{Role:} You are a professional evaluation expert.

\noindent\textbf{Task:} Given an image and a text description, evaluate whether the image faithfully reflects the content described in the text.

\noindent\textbf{Criteria:}
\begin{itemize}[leftmargin=1.5em,itemsep=1pt,topsep=2pt]
  \item \textbf{Relevance}: Does the image match the text description?
  \item \textbf{Quality}: Is the image free of artifacts and distortions?
  \item \textbf{Consistency}: Are all elements spatially coherent?
\end{itemize}

\noindent\textbf{Output:} Provide a score from 0 to 10 with a brief justification.
\end{promptbox}

\subsection{Structured Output Prompt}

% --- Example: prompt with JSON output format ---
\begin{promptbox}[System Prompt with JSON Output]
\noindent You are given a research paper abstract. Extract the following metadata and return it as a JSON object.

\vspace{0.5em}
\noindent\textbf{Expected output format:}
\begin{lstlisting}[style=json]
{
  "title": "Paper title",
  "authors": ["Author One", "Author Two"],
  "keywords": ["keyword1", "keyword2", "keyword3"],
  "task_type": "classification | generation | detection",
  "datasets": [
    {"name": "ImageNet", "split": "val", "size": 50000}
  ],
  "metrics": {
    "accuracy": 82.1,
    "f1_score": 0.79
  }
}
\end{lstlisting}

\noindent Return \textbf{only} the JSON object, with no additional text.
\end{promptbox}

% ============================================================
%  Appendix C: Theoretical Analysis (optional)
% ============================================================
\section{Theoretical Analysis}
\label{sec:theoretical_analysis}

% This section is optional. Use it when your work involves theoretical contributions.
% The template supports colored theorem boxes: Theorem (blue), Proposition (green),
% Lemma (blue), Corollary (green), Definition (yellow), Assumption (yellow), Remark (gray).
%
% Usage:  \begin{theorem}{Title}{label}  ...  \end{theorem}
% Refer:  \cref{thm:label}, \cref{prop:label}, \cref{lem:label}, etc.

\begin{theorem}{Example Theorem}{example}
Let $f: \mathcal{X} \rightarrow \mathbb{R}$ be an $L$-Lipschitz continuous function. Then for any $x, y \in \mathcal{X}$:
\begin{equation}
  |f(x) - f(y)| \leq L \|x - y\|.
\end{equation}
\end{theorem}

\begin{proof}
Your proof goes here. We omit the details for brevity.
\end{proof}
